% Einheit 4 von 5 – Umkehraufgaben (bank/binomialverteilung/e4.jsonl, zusammenbau v0.1)
%% TODO Zeitmarke Einheit 4: Marken-Zeile nicht lesbar
%% TODO Zweigzeile Teil 1: was hier gelernt wird (Satz in Schülersprache aus dem Einheitstitel) – Einheit 4
\einheitenkopf[e4]{Einheit 4 von 5 · Umkehraufgaben}
\zweigzeile{Abitur GK}
\setcounter{aufgabe}{18}

%% TODO Titel: Ich-kann-Satz fehlt in der Bank (Platzhalter: Kettenname) – Nr. 19
%% TODO Anweisung: ein Satz über der ersten Teilaufgabe fehlt in der Bank – Nr. 19
\begin{aufgabe}{Umkehraufgaben}
%% TODO \rechenplatz{n}: Schrittzahl des Grundfalls fehlt in der Bank (nur bei fester Schreibform, 2.3 b)
\begin{teile}
\teil 15\,\% der Lose gewinnen. Wie viele Lose muss man mindestens kaufen, damit mit mindestens 90\,\% Wahrscheinlichkeit mindestens ein Gewinn dabei ist? Was ist gesucht, welcher Weg passt? \kreuz{eine Wahrscheinlichkeit} \\ \kreuz{n über den Logarithmus} \\ \kreuz{p über die Wurzel} \\ \kreuz{Probieren mit Nachbarwerten}
\teil 15\,\% der Lose einer Tombola gewinnen. Stelle den Ansatz auf: Wie viele Lose muss man kaufen, damit mit mindestens 90\,\% Wahrscheinlichkeit mindestens ein Gewinn dabei ist?
\teil 64\,\% der Erwachsenen eines Landes gehen gern ins Kino, die übrigen 36\,\% nicht. Wie viele Personen müssten mindestens befragt werden, um mit einer Mindestwahrscheinlichkeit von 99\,\% wenigstens eine zu finden, die nicht gern ins Kino geht? \hfill (Abitur 2017 LK) \\ n = \leerfeld
\teil Die Wahrscheinlichkeit, dass unter 30 zufällig ausgewählten Flaschen einer Abfüllanlage keine undicht ist, soll mindestens 20\,\% betragen. Wie groß darf der Anteil undichter Flaschen höchstens sein? \hfill (Abitur 2018 GK) \\ \leerfeld
\teil X ist binomialverteilt mit 60 Versuchen und der Trefferwahrscheinlichkeit $0{,}35$. Bestimme die kleinste Zahl k mit $P(X \le k) \ge 0{,}95$ und gib beide Nachbarwerte an. k = \leerfeld
\teil X ist binomialverteilt mit 100 Versuchen und der Trefferwahrscheinlichkeit $0{,}4$. Bestimme den kleinsten Radius r, für den X mit mindestens 90\,\% Wahrscheinlichkeit höchstens um r vom Erwartungswert abweicht. r = \leerfeld
\teil 75\,\% der Erwachsenen einer Stadt besitzen ein Fahrrad. Wie viele Erwachsene muss man mindestens auswählen, damit mit mindestens 90\,\% Wahrscheinlichkeit mehr als 120 davon ein Fahrrad besitzen? \hfill (Abitur 2019 GK) \\ n = \leerfeld
\teil X ist binomialverteilt mit 50 Versuchen. Bestimme auf ganze Prozent die größte Trefferwahrscheinlichkeit, für die $P(X \le 3) \ge 0{,}5$ gilt. \leerfeld
\teil Tim behauptet: „Je mehr Lose ich kaufe, desto wahrscheinlicher ist mindestens ein Gewinn dabei.“ Beurteile die Aussage allgemein.
\teil X ist binomialverteilt; $P(X = 1)$ ist neunmal so groß wie $P(X = 0)$, und der Erwartungswert von X ist 6. Bestimme ohne Rechner die Anzahl der Versuche n und die Trefferwahrscheinlichkeit p \hfill (Abitur 2025 LK).
\end{teile}
\end{aufgabe}

%% TODO Titel: Ich-kann-Satz fehlt in der Bank (Platzhalter: Kettenname) – Nr. 20
%% TODO Anweisung: ein Satz über der ersten Teilaufgabe fehlt in der Bank – Nr. 20
\begin{aufgabe}{Trefferzahlen mit einer Einzelwahrscheinlichkeit über einer Schranke mit dem Rechner ermitteln}
\begin{teile}
\teil Ein Glücksrad bringt bei jeder Runde mit der Wahrscheinlichkeit $\frac{1}{3}$ einen Gewinn. Es werden 40 Runden gespielt; Y ist die Anzahl der Gewinnrunden. Jemand behauptet, für genau fünf Werte i sei $P(Y = i)$ größer als 10\,\%. Prüfe die Behauptung \hfill (Abitur 2017 GK).
\end{teile}
\end{aufgabe}

%% TODO Titel: Ich-kann-Satz fehlt in der Bank (Platzhalter: Kettenname) – Nr. 21
%% TODO Anweisung: ein Satz über der ersten Teilaufgabe fehlt in der Bank – Nr. 21
\begin{aufgabe}{Umkehraufgaben – Fehler finden, Begründen, Anwendung}
\begin{teile}
\teil Für „mindestens ein Gewinn mit mindestens 90\,\%“ bei 15\,\% Gewinnlosen rechnet Ole: \rechnung{0{,}85^n &\le 0{,}1 \\ n &\le \frac{\mathrm{ln}\, 0{,}1}{\mathrm{ln}\, 0{,}85} \approx 14{,}17} Er folgert, dass höchstens 14 Lose genügen. Finde den Fehler.
\teil Begründe, warum die Wahrscheinlichkeit für „kein Treffer“ mit jedem weiteren Versuch kleiner wird.
\teil Ein Rauchmelder schlägt bei Rauch mit der Wahrscheinlichkeit $0{,}9$ an, unabhängig von anderen Meldern. Wie viele Melder braucht ein Raum, damit bei Rauch mit mindestens 99,9\,\% Wahrscheinlichkeit mindestens einer anschlägt?
\end{teile}
\end{aufgabe}
